\chapter{Introduction générale}

Dans le cadre de \TODO{stage de fin d'année / PFA, nom du diplôme visé, durée exacte du
stage}, ce rapport présente le travail réalisé au sein de la Fondation OCP autour du
projet~: \emph{« Conception d'une Plateforme de Gestion des Bénéficiaires de l'Axe
Économie Sociale et Solidaire au sein de la Fondation OCP en intégrant des Exigences de
Cybersécurité »}.

La Fondation OCP pilote un ensemble de programmes sociaux et économiques regroupés sous
l'Axe Économie Sociale et Solidaire, au bénéfice de coopératives et de bénéficiaires
directs et indirects à travers le Maroc. Avant de développer une plateforme numérique de
gestion pour cet axe, il est indispensable de cadrer précisément les besoins métier, les
contraintes techniques et — étant donné la sensibilité des données traitées — les
exigences de cybersécurité et de protection des données personnelles.

C'est ce besoin de cadrage rigoureux qui a motivé la réalisation de
\textbf{Secure Requirements Studio (SRS)}~: un outil interne d'ingénierie des exigences
qui mène un entretien structuré avec les parties prenantes (en mode assisté par IA ou en
questionnaire classique hors ligne) et transforme progressivement leurs réponses en un
Cahier des Charges, une conception applicative, un diagramme de conception MVP et un
ensemble d'exigences de cybersécurité — chaque élément généré restant traçable à une
réponse effectivement donnée pendant l'entretien, plutôt qu'à un modèle générique
recopié d'un projet à l'autre.

\section{Objectifs du stage}
\TODO{Lister les objectifs pédagogiques et professionnels assignés par l'encadrant.}

\section{Structure du rapport}
Ce rapport est structuré comme suit~: le chapitre~2 présente l'organisme d'accueil ; le
chapitre~3 détaille le contexte et la problématique ; le chapitre~4 couvre l'analyse et
la conception du système ; le chapitre~5 décrit la réalisation de Secure Requirements
Studio ; le chapitre~6 explique le rôle de l'assistance par intelligence artificielle ;
le chapitre~7 détaille les aspects de cybersécurité ; le chapitre~8 présente les tests
effectués ; le chapitre~9 discute les résultats obtenus et les limites ; le chapitre~10
conclut le rapport et ouvre sur les perspectives.
